\usepackage{fontspec}
\usepackage{mathtools}
\usepackage{unicode-math}
\usepackage{csquotes}
\usepackage{accsupp}

\usepackage[lining]{FiraSans}

% do not like the lambda
\setmonofont{FiraCode}[
   Path = data/fonts/,
   Scale = 1,
   Extension = .ttf,
   Contextuals=Alternate,
   BoldFont={*-Bold},
   UprightFont={*-Regular},
]

\setmathfont{FiraMath-Regular}[
   Path = data/fonts/,
   Scale = 1.05,
   Extension = .otf
]

% TODO: find better/matching font? fira math is weird for greek stuff
\setmathfont[range={it/{greek,Greek},bfit/{greek,Greek}}]{lato}
\setoperatorfont\mathsf

\makeatletter

\def\todo#1{\textbf{\color{red}TODO:~\itshape#1}}
\def\Y{}
\def\I{}\def\S{}\def\K{}
\usepackage{xspace}

\def\InstallShortcuts{%
%
\def\textlambda{\lambda}%
\renewcommand\l[1][@main]{\ensuremath{{\color{##1}\symbf\textlambda}}}%
\def\L{\texorpdfstring{\ensuremath{\textlambda}}{\text{Lambda}}}%
\renewcommand\Y[1][@main]{\ensuremath{{\color{##1}\mathbf{Y}}}\xspace}%
\renewcommand\I[1][uulmaccent]{\ensuremath{{\color{##1}\mathbf{I}}}\xspace}%
\renewcommand\S[1][uulmaccent]{\ensuremath{{\color{##1}\mathbf{S}}}\xspace}%
\renewcommand\K[1][uulmaccent]{\ensuremath{{\color{##1}\mathbf{K}}}\xspace}%
\def\ifelse{\textsf{\textsl{ifelse}}\xspace}%
\def\pred{\textsf{\textsl{pred}}\xspace}%
\def\isZero{\textsf{\textsl{isZero}}\xspace}%
\def\T##1{\ifmmode\text{\texttt{##1}}\else\texttt{##1}\fi}%
\def\say##1{\enquote{##1}}%
\def\RED##1##2{\smash{\ensuremath{{}##1_{\rlap{\raisebox{3pt}{$\scriptstyle\color{uulmaccent}##2$}}}\mkern-.705mu{}}}}%
\def\IRED##1##2{\smash{\ensuremath{~~\;{}##1_{\raisebox{3pt}{$\scriptstyle\color{uulmaccent}##2$}}}~~}}%
\def\alphaconv{\RED=\alpha}%
\def\etaconv{\RED=\eta}%
\def\ialphaconv{\IRED=\alpha}%
\def\ibaseconv{\IRED={\phantom{\alpha}}}%
\def\ietaconv{\IRED=\eta}%
\def\betared##1{\RED={\beta_{\raisebox{-3pt}{\(\scriptscriptstyle##1\)}}}}%
\def\ibetared##1{\IRED={\beta_{\raisebox{-3pt}{\(\scriptscriptstyle##1\)}}}}%
\def\b[##1->##2]{\ensuremath{\T{{\color{uulmaccent}[##1\(\;\rightarrow\;\)##2]}}}}%
\def\@noxspace{}%
\xspaceaddexceptions{\@noxspace}%
\def\BR[##1->##2]{\ensuremath{\T{{\color{uulmaccent}[##1\@noxspace$\;\rightarrow\;$##2\@noxspace]}}}}%
% TODO: beta short red
}
\InstallShortcuts

\usepackage{xcolor}
% colors
\definecolor{uulmlogoblue}{HTML}{89a2b3}
\definecolor{uulmblue}{HTML}{7D9AAA}
\definecolor{uulmaccent}{HTML}{A9A28D}
\definecolor{green}{HTML}{56AA1C}
\definecolor{red}{HTML}{A32638}
\definecolor{orange}{HTML}{DF6D07}
\definecolor{blue}{HTML}{26547C}
\definecolor{accent}{HTML}{563C5C}

\colorlet{@main}{accent}

\preto\input@path{{lib/beamer-themes}{lib/tikzpingus/tex}{lib/fancyqr}{lib/algorithms}}
\usepackage{lib/fancyqr/fancyqr}
\usepackage{lib/algorithms/nice-algo}
\SetAlgoVlined
\RestyleAlgo{plain}

%---required packages & variable definitions------------------------------------
\usepackage{forest}
\usepackage{xcolor}
\usetikzlibrary{angles}
\definecolor{drawColor}{RGB}{64, 64, 64}
\newcommand{\circleSize}{0.165em}
\newcommand{\angleSize}{0.8em}
%-------------------------------------------------------------------------------
%---Define the style of the tree------------------------------------------------
\forestset{
	/tikz/mandatory/.style={
		circle,fill=drawColor,
		draw=drawColor,
		inner sep=\circleSize
	},
	/tikz/optional/.style={
		circle,
		fill=white,
		draw=drawColor,
		inner sep=\circleSize
	},
	featureDiagram/.style={
		for tree={
			text depth = 0,
			parent anchor = south,
			child anchor = north,
			draw = drawColor,
			rounded corners=2pt,
			edge = {draw=drawColor},
		}
	},
	/tikz/abstract/.style={
		fill = @main!15!white,
		draw = drawColor
	},
	/tikz/concrete/.style={
		fill = none,
		draw = drawColor
	},
	mandatory/.style={
		edge label={node [mandatory] {} }
	},
	optional/.style={
		edge label={node [optional] {} }
	},
	or/.style={
		tikz+={
			\path (.parent) coordinate (A) -- (!u.children) coordinate (B) -- (!ul.parent) coordinate (C) pic[fill=drawColor, angle radius=\angleSize]{angle};
		}
	},
	/tikz/or/.style={
	},
	alternative/.style={
		tikz+={
			\path (.parent) coordinate (A) -- (!u.children) coordinate (B) -- (!ul.parent) coordinate (C) pic[draw=drawColor, angle radius=\angleSize]{angle};
		}
	},
	/tikz/alternative/.style={
	},
	/tikz/placeholder/.style={
	},
	collapsed/.style={
		rounded corners,
		no edge,
		for tree={
			fill opacity=0,
			draw opacity=0,
			l = 0em,
		}
	},
	/tikz/hiddenNodes/.style={
		midway,
		rounded corners,
		draw=drawColor,
		fill=white,
		minimum size = 1.2em,
		minimum width = 0.8em,
		scale=0.9
	},
}
%-------------------------------------------------------------------------------

\def\pingulibpathprefix{lib/tikzpingus/tex/}
\usepackage[glows]{lib/tikzpingus/tex/tikzpingus}
\makeatletter
\usetikzlibrary{tikzmark,arrows.meta}
% The aux file is read inside the document group, where every \csname of a
% new mark creates a local \relax entry. With tens of thousands of
% remember-picture marks this exhausts TeX's save stack (save size), so we
% scope each entry; the actual definitions are global anyway.
\let\fp@savepicturepage\savepicturepage
\def\savepicturepage#1#2{\begingroup\fp@savepicturepage{#1}{#2}\endgroup}
\let\fp@pgfsyspdfmark\pgfsyspdfmark
\def\pgfsyspdfmark#1#2#3{\begingroup\fp@pgfsyspdfmark{#1}{#2}{#3}\endgroup}
\tikzset{
   @O/.style={overlay,remember picture},
   T/.style={uulmaccent!60!lightgray,font=\footnotesize\itshape\sffamily,line cap=round},
   lcr/.style={line cap=round}
}
\def\HStrut{\vphantom{\{tg}}% }<- wrong hl
\def\Snode#1{\tikzmarknode{#1}\HStrut}%
\def\TSnode#1{\subnode{#1}\HStrut}%
\def\bnode#1#2{\rnode{#1}{\hsGetStyle{basic}{#2}}}%
\def\rnode#1#2{\Snode{#1}#2\Snode{#1@}}
\def\TBM#1#2{\Snode{#1}#2\Snode{#1@}}
\def\hlcolor{@main}\def\hlopa{0.25}\def\hlhs{.4mm}% horizonal shift
\newcommand\hlrawbehindcodeunder[4][0pt]{\setbox0=\hbox{\pgfinterruptpicture\HStrut\endpgfinterruptpicture}\fill[\hlcolor,opacity=\hlopa,rounded corners=1pt] ([shift={(-\hlhs,-#1-1.6pt)}]#4#2) rectangle ([shift={(\hlhs,-#1+1.6pt)}]#4#2@); \coordinate (#3) at (#4#2@);
\path ([shift={(-\hlhs,-#1+.95\dp0)}]#4#2) -- ([shift={(\hlhs,-#1+.95\ht0)}]#4#2@) coordinate[pos=.5,yshift=1mm] (#3_up);
\path ([shift={(-\hlhs,-#1-.95\dp0)}]#4#2) -- ([shift={(\hlhs,-#1+.95\ht0)}]#4#2@) coordinate[pos=.5,yshift=-1mm-1.6pt] (#3_down); }
\newcommand\hlbehindcodeunder[3][0pt]{\hlrawbehindcodeunder[#1]{#2}{#3}{pic cs:}}
\newcounter{ULcounter}
\def\UL#1{\Snode{@ul@\theULcounter}#1\Snode{@ul@\theULcounter @}\tikz[@O]{\hlbehindcodeunder{@ul@\theULcounter}{@>}}\stepcounter{ULcounter}}

\def\G#1{\textcolor{gray}{#1}}
\def\Gopen{\textcolor{gray}{(}}
\def\Gclose{\textcolor{gray}{)}}
\newcommand\hlbehindcode[3][2pt]{\setbox0=\hbox{\pgfinterruptpicture\HStrut\endpgfinterruptpicture}\fill[\hlcolor,opacity=\hlopa,rounded corners=1pt] ([shift={(-\hlhs,-#1-.95\dp0)}]pic cs:#2) rectangle ([shift={(\hlhs,#1+.95\ht0)}]pic cs:#2@); \coordinate (#3) at ([yshift=.475\ht0]pic cs:#2@);
\path ([shift={(-\hlhs,#1+.95\dp0)}]pic cs:#2) -- ([shift={(\hlhs,#1+.95\ht0)}]pic cs:#2@) coordinate[pos=.5,yshift=1mm] (#3_up);\path ([shift={(-\hlhs,-#1-.95\dp0)}]pic cs:#2) -- ([shift={(\hlhs,-#1+.95\ht0)}]pic cs:#2@) coordinate[pos=.5,yshift=-1mm] (#3_down); }% hardcoded 1mm padding
% just underlines as there are no vert buffers:D

\newsavebox\pingubox


\let\hl\relax
\usepackage{lib/latex-underline/fancy-underline}

% Konfiguration der Unterstriche bei Links
\def\fancyul@depth{.965\p@}
\def\fancyul@contour{.2\p@}
\colorlet{tcbcolback}{white}
\def\fancyulbackground{tcbcolback}% so that it adapts with the beamer
% we must hook the background for the boxes

\colorlet{url@line@color}{uulmaccent}
\fancyulcolor{url@line@color!45!\fancyulbackground}
\newdimen\ul@surround \ul@surround=1.35\p@
% we make the auto detection of the depth bigger
\def\update@ul{\setbox\z@\hbox{{(j}}\edef\fancyul@thickness{\the\dimexpr.2\dp\z@}}
\def\rawlink#1{\leavevmode\kern-\ul@surround{\update@ul\kern\ul@surround\relax\fancyul{#1}}}
% define a link macro that can be used with href and fancyul
\DeclareRobustCommand*\link{\hyper@normalise\link@do}
\def\link@do#1#2{\leavevmode\kern-\ul@surround\href{#1}{\update@ul\kern\ul@surround\relax\fancyul{#2}}}% thinspace for some nicer padding
\def\xlink@color{hs@colors@lst@keywordA}
\DeclareRobustCommand*\Xlink{\hyper@normalise\Xlink@do}
\def\Xlink@do#1#2{\leavevmode\kern-\ul@surround\href{#1}{\update@ul\kern\ul@surround\relax\color{\xlink@color}\fancyul{#2}}}% thinspace for some nicer padding
\DeclareRobustCommand*\hlink{\hyper@normalise\hlink@do}
\def\hlink@do#1#2{\leavevmode\kern-\ul@surround\hyperlink{#1}{\update@ul\kern\ul@surround\relax\fancyul{#2}}}%
\usepackage{letltxmacro}
\AtBeginDocument{%
\LetLtxMacro\basicurl\url
% \def\url{\begingroup\catcode`{11\catcode`}11\obeyspaces\obeylines\url@do}%
% \def\url@do#1{\endgroup\link{#1}{#1}}%
\DeclareRobustCommand*\url{\hyper@normalise\url@do}
\def\url@do#1{\leavevmode\kern-\ul@surround\href{#1}{\update@ul\kern\ul@surround\relax\fancyul{#1}}}%
\long\def\beamer@ref#1{\hlink@do{#1}{{\beamer@origref*{#1}}}}% the extra set of braces is for soul so that it can work with braces
}
\def\abs#1{\left|#1\right|}
\def\N{\ensuremath{\mathbb{N}}}

\colorlet{@thegray}{uulmaccent!20!gray}
\def\info#1{\textcolor{@thegray}{\scriptsize\colorlet{@}{@thegray}#1}}
\def\HL#1#2{\T{\link{#1}{#2}}}
\def\DescHl{}
\def\Mix{15}\def\MixB{5}

\def\DescHeight{1.75cm}\def\DescText#1#2{\textbf{\vphantom{A}\smash{\DescHl\TLnk{#1}{#1}}}\nodepart{second}\parbox[c][\DescHeight]{.9\linewidth}{\footnotesize\DescHl\centering\T{#2}\endgraf}}
\def\RawDescText#1#2{\textbf{\vphantom{A}\smash{\DescHl#1}}\nodepart{second}\parbox[c][\DescHeight]{.9\linewidth}{\footnotesize\DescHl\centering\T{#2}\endgraf}}
\tikzset{@/.style={}}%
\def\BSep{5mm}
\def\RSep{9mm}

\def\EnumPhrase{\link{https://stackoverflow.com/questions/58761160/why-isnt-enum-typeclass-a-subclass-of-ord-typeclass}{Enum is no subclass of Ord?}}
\def\Yields{~~\textcolor{uulmaccent}{\faCaretRight}~~}%
% \gappto\OuterStuff{\onslide<#1|handout:#2>{#4}}
\def\ActiveOn#1#2#3{{\def\Mix{15}\def\MixB{5}\def\DescHl{}\only<3->{\def\Mix{5}\def\MixB{0}\def\DescHl{\color{gray}}}\only<#1|handout:#2>{\tikzset{K/.append style={black}}\def\Mix{15}\def\MixB{5}\def\DescHl{}}#3}}
\newcommand\Lnk[3][Prelude]{\T{\href{https://hackage.haskell.org/package/base-4.21.0.0/docs/#1.html\#v:#2}{#3}}}
\newcommand\LLnk[2][Prelude]{\T{\href{https://hackage.haskell.org/package/base-4.21.0.0/docs/#1.html\#v:#2}{\fancyul{#2}}}}
\newcommand\TLnk[3][Prelude]{\T{\href{https://hackage.haskell.org/package/base-4.21.0.0/docs/#1.html\#t:#2}{#3}}}
\newcommand\MLnk[2][Prelude]{\T{\href{https://hackage.haskell.org/package/base-4.21.0.0/docs/#1.html}{\fancyul{#2}}}}
\newcommand\TTLnk[2][Prelude]{\T{\href{https://hackage.haskell.org/package/base-4.21.0.0/docs/#1.html\#t:#2}{\fancyul{#2}}}}
\def\BLnk#1#2{\T{\href{https://hackage.haskell.org/package/base-4.21.0.0/docs/Control-Applicative.html\#v:#1}{#2}}}
\tikzset{K/.style={rectangle,draw,rounded corners=3.5pt,text width=3cm,minimum height=1.9cm,align=center,inner sep=5pt,rectangle split,rectangle split parts=2,rectangle split part fill={#1!\Mix!white,#1!\MixB!}}}

\def\CodeFileMarker#1{\CodeMarker{\llap{\faFileCodeO~\thinspace}#1}}
\def\CodeMarker#1{{\color{lightgray!75!gray}\sffamily\footnotesize#1}}
\colorlet{@filemarker}{lightgray!75!gray}
\let\curpath\@empty

\def\StartExcursion{%
   \begingroup
   \setbeamercolor{background canvas}{bg=gray!15}%
   \appto\PlaceFramenumberUpperLeft{\smash{\rlap{\tikz[remember picture,overlay]\node[below right,font=\tiny,gray!50,scale=.75] at(current page.north west) {excursion};}}}%
}
\def\EndExcursion{\endgroup}